\pdfoutput=1
\documentclass[12pt]{article}

\usepackage{amsmath, amssymb, amsthm}
\usepackage{graphicx}
\usepackage{booktabs}
\usepackage{array}
\usepackage{longtable}
\usepackage{calc}
\usepackage{etoolbox}
\makeatletter
\patchcmd\longtable{\par}{\if@noskipsec\mbox{}\fi\par}{}{}
\makeatother
\IfFileExists{footnotehyper.sty}{\usepackage{footnotehyper}}{\usepackage{footnote}}
\makesavenoteenv{longtable}
\usepackage[colorlinks=true, linkcolor=blue, citecolor=blue, urlcolor=blue]{hyperref}
\usepackage{geometry}
\geometry{margin=1in}

\makeatletter
\newsavebox\pandoc@box
\newcommand*\pandocbounded[1]{%
  \sbox\pandoc@box{#1}%
  \Gscale@div\@tempa{\textheight}{\dimexpr\ht\pandoc@box+\dp\pandoc@box\relax}%
  \Gscale@div\@tempb{\linewidth}{\wd\pandoc@box}%
  \ifdim\@tempb\p@<\@tempa\p@\let\@tempa\@tempb\fi
  \ifdim\@tempa\p@<\p@\scalebox{\@tempa}{\usebox\pandoc@box}%
  \else\usebox{\pandoc@box}%
  \fi
}
\makeatother
\providecommand{\tightlist}{%
  \setlength{\itemsep}{0pt}\setlength{\parskip}{0pt}}

\title{Dual Entropy Gradients and the Tholonic Void-Plenum Equivalence}
\author{J. W. Milton\\[4pt]{\small Clarity Coalition}}
\date{\small 13 June 2026 \\ v1.0}

\begin{document}
\maketitle

\noindent\textbf{Keywords:} entropy; dual entropy; tholonic model; N-D-C framework; conceptual entropy; material entropy; void-plenum equivalence; animate matter; active principles; information theory; cosmology; self-organization

\begin{abstract}
The Tholonic N-D-C framework, when applied at the cosmological scale,
predicts two opposing entropy gradients operating simultaneously and
structurally. The material domain, identified with the C
(Contribution/Integration) role at the cosmic scale, undergoes entropy
increase: it disperses, differentiates, and moves from minimum entropy
(the material singularity) toward maximum entropy (maximum diffusion,
heat death). The conceptual domain, identified with the D
(Definition/Limitation) role at the cosmic scale, undergoes entropy
decrease: it organizes, becomes more specific, more coherent, and more
structured over time, moving from maximum entropy (the conceptual void:
one concept, zero differentiation, zero organization) toward minimum
entropy (maximum structure, maximum connection, maximum coherence).

This duality is not an additional postulate of the tholonic model. It is
a structural consequence of the N-D-C role assignment: C is the
dispersive, expansive, accumulating pole, and D is the constraining,
organizing, specifying pole. At the cosmic scale, these roles map
directly onto the material and conceptual domains respectively. The N
state of the universe at any moment is the balance between these two
opposing gradients.

The paper further argues that the starting point of the conceptual
expansion (the void: a zero-dimensional point of zero differentiation)
and its endpoint (infinite concepts organized into infinite mutual
connections) are topologically and informationally equivalent. Both are
states of zero Shannon entropy in different senses: the void is
maximally underdetermined (all possibilities equally probable), while
the plenum is maximally determined (all possibilities fully realized).
This equivalence is supported by Kuiper's theorem for
infinite-dimensional Hilbert spaces, the self-duality of projective
geometry in the limit, and the behavior of Shannon entropy at its two
extremes.

The paper is explicitly exploratory. The structural argument (dual
entropy gradients as a consequence of N-D-C role assignment) is derived
from the established tholonic framework. The cosmological claims
(void-plenum equivalence, the material/conceptual mapping) are argued
structurally and by mathematical analogy, not by formal proof. Open
questions identify what would be required to elevate the central claims
from structural argument to theorem.
\end{abstract}

\section{Introduction}\label{introduction}

The second law of thermodynamics states that entropy in an isolated
system does not decrease. This is among the best-confirmed principles in
physics. It is also, on inspection, a statement about a specific domain:
the material domain, specifically the domain of energy distribution in
physical systems.

The question this paper addresses is whether the same principle, or its
converse, applies to a different domain: the domain of concepts, ideas,
and information. If the material domain is characterized by entropy
increase, is the conceptual domain characterized by entropy decrease?
And if so, is this duality a consequence of some deeper structural
principle, or is it an independent empirical observation?

The Tholonic N-D-C framework
\href{https://github.com/baardev/truevalue/raw/main/docnav/Research/papers/3_minimal-recursive-triadic-framework/3_minimal-recursive-triadic-framework.pdf}{Mil24c}
provides a structural answer. The framework assigns three irreducible
functional roles to any self-organizing recursive system: N
(Negotiation: the emergent stable state), D (Definition/Limitation: the
constraining, organizing, specifying role), and C
(Contribution/Integration: the expansive, accumulating, dispersive
role). At the cosmic scale, the material domain exhibits the properties
of the C role: it disperses, expands, and differentiates away from an
initial concentration. The conceptual domain exhibits the properties of
the D role: it constrains, organizes, and becomes more specifically
structured over time. The dual entropy gradients are therefore not an
independent observation but a consequence of role assignment.

This paper also argues that the two endpoints of the conceptual entropy
trajectory are equivalent to each other and to the single starting point
of the material trajectory. The material universe begins as a point
singularity (minimum volume, minimum material entropy) and expands
toward maximum material entropy. The conceptual universe begins as the
void (zero concepts, maximum conceptual entropy) and organizes toward
minimum conceptual entropy. The endpoint of full conceptual organization
(infinite concepts connected to infinite concepts) is, by several
mathematical arguments, equivalent to the starting void. The tholonic
recursive structure both predicts and explains this equivalence: at the
largest scale, the child N of the entire hierarchy must be equivalent to
the parent N that generated it.

The connections to Newton's ``active principles'' {[}New17{]} and to
Sidis's animate/inanimate distinction {[}Sid25{]} are treated in
\href{https://github.com/baardev/truevalue/raw/main/docnav/Research/papers/12_newton-tholonic-framework/12_newton-tholonic-framework.pdf}{Paper
12}
\href{https://github.com/baardev/truevalue/raw/main/docnav/Research/papers/12_newton-tholonic-framework/12_newton-tholonic-framework.pdf}{Mil26a}.
The present paper takes the structural argument further and examines its
cosmological implications independently.

\textbf{What this paper provides.} A structural derivation of dual
entropy gradients from the N-D-C role assignment. A mathematical
argument for the void-plenum equivalence via Kuiper's theorem, Shannon
entropy, and projective duality. A set of tholonic predictions about the
relationship between material and conceptual domains. Explicit open
questions separating what is argued from what remains to be proved.

\textbf{What this paper does not provide.} A formal proof of void-plenum
equivalence. A physical model of how the conceptual domain interacts
with the material domain. A cosmological model making specific
quantitative predictions about entropy values or timescales. The paper
is structural and exploratory.

\textbf{Organization.} Section 2 reviews the N-D-C role assignment and
the cosmic-scale mapping. Section 3 argues the case for material entropy
increase as C expansion. Section 4 argues the case for conceptual
entropy decrease as D organization. Section 5 examines Newton's active
principles and Sidis's animate matter as the interface between the two
domains. Section 6 develops the void-plenum equivalence argument.
Section 7 provides the mathematical grounding. Section 8 states tholonic
predictions and their status. Section 9 poses open questions. Section 10
concludes.

\begin{center}\rule{0.5\linewidth}{0.5pt}\end{center}

\section{The N-D-C Role Assignment at the Cosmic Scale}\label{the-n-d-c-role-assignment-at-the-cosmic-scale}

The tholonic triad assigns three functional roles to any recursive
self-organizing system:

\begin{itemize}
\tightlist
\item
  \(N\) (Negotiation): the emergent stable state. It is not directly
  measured; it arises from the interaction of D and C. It is
  simultaneously the product of one level and the source of the next.
\item
  \(D\) (Definition/Limitation): the constraining, bounding, specifying
  role. D governs what a system IS: its boundaries, identity, and
  organization. D reduces possibility to actuality by specification.
\item
  \(C\) (Contribution/Integration): the accumulating, expansive,
  dispersive role. C governs what a system DOES: its outputs,
  relationships, and extensions. C increases the scope of actualized
  states.
\end{itemize}

The fundamental stability principle is \(D \approx C\): systems are most
stable when their constraining and expansive poles are balanced.
Imbalance in either direction increases energy cost and degrades the N
state.

At the cosmic scale, the assignment is:

\begin{longtable}[]{@{}
  >{\raggedright\arraybackslash}p{(\linewidth - 4\tabcolsep) * \real{0.3333}}
  >{\raggedright\arraybackslash}p{(\linewidth - 4\tabcolsep) * \real{0.3333}}
  >{\raggedright\arraybackslash}p{(\linewidth - 4\tabcolsep) * \real{0.3333}}@{}}
\toprule\noalign{}
\begin{minipage}[b]{\linewidth}\raggedright
Role
\end{minipage} & \begin{minipage}[b]{\linewidth}\raggedright
Cosmic domain
\end{minipage} & \begin{minipage}[b]{\linewidth}\raggedright
Entropy behavior
\end{minipage} \\
\midrule\noalign{}
\endhead
\bottomrule\noalign{}
\endlastfoot
\(C\) & Material domain: energy, matter, spacetime & Increases
(disperses from singularity) \\
\(D\) & Conceptual domain: information, structure, relations & Decreases
(organizes from void) \\
\(N\) & The universe as self-organizing system & The balance between the
two gradients \\
\end{longtable}

This assignment is not arbitrary. The C role is characterized by
accumulation and extension: it grows the scope of what exists. Material
entropy increase is precisely this: the universe accumulates more and
more realized microstates, extending its physical configuration space.
The D role is characterized by specification and constraint: it reduces
the space of possibilities to a more precisely organized subset.
Conceptual entropy decrease is precisely this: the space of concepts
becomes more organized, more internally constrained, and more
specifically structured over time.

\begin{figure}
\centering
\pandocbounded{\includegraphics[width=\linewidth,keepaspectratio]{figures/13_tholonic-triad-cosmic.png}}
\caption{Figure 1: The Tholonic Triad at Cosmic Scale. N (Universe, top
vertex) is the emergent balance of D (Conceptual Domain, bottom-left)
and C (Material Domain, bottom-right). The bottom edge marks the animate
matter interface where the two entropy gradients meet.}
\end{figure}

\emph{Figure 1: The tholonic triad at the cosmic scale. The D and C
roles map directly to the conceptual and material domains. The bottom
edge, where D and C meet, corresponds to animate matter: the interface
at which D-domain entropy decrease achieves material instantiation.}

\begin{center}\rule{0.5\linewidth}{0.5pt}\end{center}

\section{Material Entropy Increase as C Expansion}\label{material-entropy-increase-as-c-expansion}

The material universe began, in the standard cosmological model, as a
state of extreme concentration: high energy density, minimum spatial
extension, minimum material entropy (in the sense of Boltzmann: few
accessible microstates). Since the initial singularity, the universe has
expanded. Energy has dispersed. Structures have formed and dissolved.
The number of accessible microstates has grown. Material entropy has
increased and continues to increase.

In tholonic terms, this is the C role operating at the cosmic scale. The
C pole of the tholonic triad is the accumulating, extending,
contributing pole. It disperses from a concentrated source, increases
the scope of what exists, and adds to the running state. The material
Big Bang is the initial condition of the C expansion: a maximally
concentrated C source beginning its accumulation outward.

The trajectory is:

\[\text{Material singularity (C}_\text{min}\text{)} \longrightarrow \text{Expansion} \longrightarrow \text{Heat death (C}_\text{max}\text{)}\]

At heat death, material entropy is maximized: all energy is uniformly
distributed, no free energy remains, no further differentiation is
possible. The C expansion is complete. Every possible material
microstate has been realized and dispersed.

This trajectory is monotonic and irreversible in the C direction. There
is no material return to the singularity. The tholonic balance principle
predicts this: once C has expanded beyond the \(D \approx C\) balance
point, the system cannot recover the N state without a D input of equal
magnitude. In a universe where D is the conceptual domain
(non-material), that input cannot act on the material domain directly
except through the interface of animate/vegetative systems.

\begin{center}\rule{0.5\linewidth}{0.5pt}\end{center}

\section{Conceptual Entropy Decrease as D Organization}\label{conceptual-entropy-decrease-as-d-organization}

The conceptual domain begins from a different starting point. The
simplest possible concept is the void: not the physical vacuum (which
has structure, energy, and quantum fluctuations) but the pure conceptual
void: zero differentiation, zero relationship, zero organization. Every
concept is equally probable; no concept is distinguished from any other.
This is maximum conceptual entropy in the Shannon sense: maximum
uncertainty, minimum information, a uniform distribution over all
conceptual possibilities.

From this starting point, the conceptual trajectory is the inverse of
the material one:

\[\text{Conceptual void (D}_\text{max entropy}\text{)} \longrightarrow \text{Organization} \longrightarrow \text{Plenum (D}_\text{min entropy}\text{)}\]

Each act of conceptual organization reduces conceptual entropy: it
distinguishes one concept from another, establishes a relationship,
specifies a constraint. Mathematics, physics, logic, and language are
all instruments of this process. They do not discover pre-existing
structure passively; they participate in the D-domain process of
progressive specification.

The trajectory is driven by the D role's defining property:
specification. D reduces the space of possibilities to a more precisely
organized subset. Every theorem proved, every relationship identified,
every distinction drawn is a decrease in conceptual entropy. The D
expansion is toward minimum conceptual entropy: a state in which every
concept is fully specified by its relationships to every other concept.

This trajectory is also monotonic: conceptual entropy decreases over
time, in the sense that the organized body of concepts only grows more
organized. Individual concepts can be discarded or revised, but the
overall degree of organization of the conceptual space increases. What
is lost in one branch of inquiry is typically recovered and exceeded in
another.

\begin{figure}
\centering
\pandocbounded{\includegraphics[width=\linewidth,keepaspectratio]{figures/13_dual-entropy-trajectories.png}}
\caption{Figure 2: Dual Entropy Trajectories. The material (C) domain
begins at minimum entropy and rises toward heat death (orange). The
conceptual (D) domain begins at maximum entropy and falls toward the
plenum (blue). They cross at the N state, the region where animate
matter and active principles operate.}
\end{figure}

\emph{Figure 2: The two opposing entropy trajectories. Material entropy
(orange, C domain) increases monotonically from the singularity toward
heat death. Conceptual entropy (blue, D domain) decreases monotonically
from the void toward the plenum. The crossing point is the N state: the
region of local entropy reversal occupied by animate matter and Newton's
active principles.}

\begin{center}\rule{0.5\linewidth}{0.5pt}\end{center}

\section{Newton's Active Principles and Sidis's Animate Matter}\label{newtons-active-principles-and-sidiss-animate-matter}

Newton's unresolved question, posed in Query 31 of the Opticks
{[}New17{]}, was the mechanism of ``active principles'': non-mechanical
organizing forces that produce coherence, structure, and persistence
that mechanical operations cannot generate. In the tholonic reading,
active principles are N states. They are the local instantiations of
D-domain organization acting within the C-domain material world. A
living organism, a crystal, a stable molecular structure: each is a site
where the D-domain entropy decrease has achieved material instantiation,
temporarily reversing local material entropy against the global C-domain
trend.

Sidis {[}Sid25{]}, approaching the same boundary from statistical
mechanics, proposed that animate matter is matter that locally reverses
entropy. His argument was that this reversal is statistically necessary
given time-symmetric microscopic laws. In the tholonic reading, the
statistical necessity argument corresponds to the structural necessity
of N emergence: given D and C interacting under the minimality
conditions of
\href{https://github.com/baardev/truevalue/raw/main/docnav/Research/papers/3_minimal-recursive-triadic-framework/3_minimal-recursive-triadic-framework.pdf}{Paper
3}
\href{https://github.com/baardev/truevalue/raw/main/docnav/Research/papers/3_minimal-recursive-triadic-framework/3_minimal-recursive-triadic-framework.pdf}{Mil24c},
N must emerge at some balance point. Animate matter is the material
trace of that emergence: the N state made locally visible in the C
domain.

The dual entropy framework extends both Newton's and Sidis's
observations: they were each identifying the interface between the two
domains. Animate matter (Newton's vegetable spirit, Sidis's
entropy-reversing systems) is not simply a local exception to the second
law. It is the point at which the D-domain entropy decrease gradient
achieves material expression, temporarily overriding the C-domain
entropy increase gradient at the local scale. The interface is the N
state.

\begin{center}\rule{0.5\linewidth}{0.5pt}\end{center}

\section{The Void-Plenum Equivalence}\label{the-void-plenum-equivalence}

The most striking structural claim of this paper is that the starting
point of the conceptual trajectory and its endpoint are equivalent. This
requires argument.

\textbf{The conceptual void} is the state of maximum conceptual entropy:
zero differentiation, zero organization, zero structure. Every possible
concept is equally likely. No concept has priority. In Shannon's terms,
this is the uniform distribution over all possibilities: maximum entropy
\(H = \log_2 N\) where \(N\) is the total number of possible concepts
(or, in the limit, \(H \to \infty\) as \(N \to \infty\)). Geometrically,
this is a zero-dimensional point: no distinctions have been drawn, so no
dimensions of conceptual space exist.

\textbf{The conceptual plenum} is the state of minimum conceptual
entropy: all concepts fully specified by their relationships to all
other concepts. Every distinction has been drawn. Every relationship has
been established. Every concept is fully determined by its position in
the network of all concepts. In Shannon's terms, this is a degenerate
distribution where every concept's state is fully predictable from the
others: entropy \(H = 0\). Geometrically, this is an
infinite-dimensional space: every possible distinction defines a
dimension.

The paradox is that both endpoints have entropy \(H = 0\) (or its
limiting analogue). The void has zero entropy because there is only one
state (undifferentiated nothingness: all possible concepts are the same
state, namely ``no concept''). The plenum has zero entropy because every
concept is fully determined by its connections: zero residual
uncertainty. Both are maximally determined, from opposite directions.

This is the tholonic prediction: the parent N at the top of the
hierarchy (the void, undifferentiated source) and the child N at the
bottom (the plenum, fully realized actuality) are the same entity viewed
from opposite ends of the recursive descent. The full tholonic cycle:
void (parent N) differentiates into D (conceptual domain) and C
(material domain), which interact to produce organized reality (child
N), which, when the hierarchy is fully realized, contains all possible N
states, and is therefore equivalent to the void that contained all
possible N states in potential form.

The distinction between the void and the plenum is the distinction
between potential and actual: both are complete, but one is complete in
latency and the other in realization. A universe that has realized all
possible N states through D-C negotiation is, from the outside,
indistinguishable from the void that contained all of them as
undifferentiated potential.

\begin{center}\rule{0.5\linewidth}{0.5pt}\end{center}

\section{Mathematical Grounding}\label{mathematical-grounding}

Three mathematical results support the void-plenum equivalence. None
constitutes a proof of the full claim, but each provides structural
evidence that the equivalence is not merely philosophical.

\subsection{Kuiper's Theorem}\label{kuipers-theorem}

Kuiper's theorem (1965) states that the unitary group \(U(H)\) of an
infinite-dimensional separable Hilbert space \(H\) is contractible in
the norm topology. This means \(U(H)\) has the homotopy type of a point:
all of its homotopy groups vanish. An infinite-dimensional Hilbert
space, under its natural unitary symmetry group, has the same
topological complexity as a single point.

More precisely: in finite dimensions, \(U(n)\) has rich homotopy
(\(\pi_1(U(n)) = \mathbb{Z}\), etc.). As \(n \to \infty\), the homotopy
groups stabilize and then vanish. The topological structure of
infinite-dimensional unitary symmetry collapses to the trivial structure
of a point. The infinite-dimensional conceptual space, under its natural
symmetry operations, is topologically a point.

This is not the same as saying an infinite-dimensional space is a point.
It says that under the natural operations defined on it, its invariant
structure is point-like. The void (a zero-dimensional point) and the
plenum (an infinite-dimensional space with collapsed homotopy) share
this invariant.

\subsection{Shannon Entropy at Extremes}\label{shannon-entropy-at-extremes}

Shannon entropy is:

\[H(p) = -\sum_{i} p_i \log_2 p_i\]

It achieves its maximum for a uniform distribution over \(n\) states:
\(H = \log_2 n\). As \(n \to \infty\), \(H \to \infty\). This is the
void in one sense: infinite uncertainty, no structure.

It achieves its minimum (\(H = 0\)) for a degenerate distribution: all
probability on one state. This is both the void in another sense (only
one possible state: undifferentiated nothingness) and the plenum (only
one possible state: fully determined by all relationships).

The symmetry is: the void is \(H = 0\) because there are no distinctions
(\(n = 1\)). The plenum is \(H = 0\) because all distinctions are fully
determined. Between them, \(H\) rises through the conceptual expansion
(as distinctions are drawn and possibilities proliferate) and then falls
back to zero as the system becomes fully organized.

\begin{figure}
\centering
\pandocbounded{\includegraphics[width=\linewidth,keepaspectratio]{figures/13_shannon-void-plenum.png}}
\caption{Figure 3: Shannon Entropy Along the Conceptual Expansion. Both
the void (left) and the plenum (right) are states of H=0, approached
from opposite directions. The peak marks maximum uncertainty:
distinctions proliferating but not yet organized. The double-headed
arrow marks the structural equivalence of the two endpoints.}
\end{figure}

\emph{Figure 3: Shannon entropy along the conceptual expansion axis.
Both the void (\(H = 0\): one undifferentiated state) and the plenum
(\(H = 0\): every state fully determined) are states of zero entropy,
approached from opposite directions. The U-shaped curve is the
trajectory of the D domain from Section 4. The void and plenum are the
two fixed points of the projective duality discussed in Section 7.3.}

\subsection{Projective Self-Duality}\label{projective-self-duality}

In projective geometry, points and hyperplanes are dual: every theorem
about points corresponds to a theorem about hyperplanes under the
duality transform. A projective point is ``a line through the origin''
and a projective hyperplane is ``a codimension-1 subspace.'' In \(n\)
dimensions, this duality is exact.

In the infinite-dimensional limit, projective duality becomes
self-referential in specific senses. The dual of an infinite-dimensional
projective space under certain function-space topologies is homeomorphic
to the original space (by the Fréchet-Riesz representation theorem for
Hilbert spaces: \(H \cong H^*\)). The space is its own dual. A fully
connected infinite-dimensional conceptual space, where every concept
determines a hyperplane of related concepts, is self-dual under the same
duality that maps a point to its polar hyperplane.

The void (a single undifferentiated point, no polar hyperplane defined)
and the plenum (every point has a polar hyperplane, the space is
self-dual) are the two fixed points of this duality operation. They are
the points at which the duality map is trivially satisfied: the void
because there is nothing to map, the plenum because the map is the
identity.

\begin{center}\rule{0.5\linewidth}{0.5pt}\end{center}

\section{Tholonic Predictions}\label{tholonic-predictions}

The structural argument produces several predictions about the
relationship between the material and conceptual domains. These are
stated here with their current status.

\textbf{Prediction 1: Animate matter marks D-C interface locations.}
Sites of local material entropy reversal (animate matter, living
systems, self-organizing structures) are sites where D-domain
organization achieves material instantiation. They should be analyzable
as local N states with measurable D-C balance. \emph{Status: testable
via the tholonic balance functional \(B(D,C)\) applied to biological and
self-organizing systems. Supported qualitatively by the vegetation
manuscript analysis
(\href{https://github.com/baardev/truevalue/raw/main/docnav/Research/papers/12_newton-tholonic-framework/12_newton-tholonic-framework.pdf}{Paper
12}
\href{https://github.com/baardev/truevalue/raw/main/docnav/Research/papers/12_newton-tholonic-framework/12_newton-tholonic-framework.pdf}{Mil26a},
Appendix B).}

\textbf{Prediction 2: The conceptual entropy of a domain decreases as a
function of the number of established relationships per concept.} As a
field of inquiry matures, the ratio of established relationships to
total concepts should increase monotonically. The conceptual entropy of
mature fields should be measurably lower than that of emerging fields.
\emph{Status: testable via bibliometric or knowledge graph analysis. Not
yet tested.}

\textbf{Prediction 3: The phi-optimal balance point (\(D \approx C\))
marks the transition between material-entropy-dominated and
conceptual-entropy-dominated dynamics.} Below the balance point,
material entropy increase dominates (C-dominant, dispersive behavior).
Above it, conceptual entropy decrease can temporarily dominate
(D-dominant, organizing behavior). The 61.8\% threshold of the tholonic
balance functional marks this transition. \emph{Status: structurally
derived from the N-D-C framework. Quantitative test requires identifying
measurable proxies for D and C in physical and biological systems.}

\textbf{Prediction 4: The two endpoints of the full tholonic hierarchy
are informationally indistinguishable.} A universe that has realized all
possible N states through D-C negotiation should be informationally
equivalent to the undifferentiated void. \emph{Status: argued via
Kuiper's theorem and Shannon entropy (Section 7). Not formally proved.
This is the most speculative prediction in the paper.}

\begin{center}\rule{0.5\linewidth}{0.5pt}\end{center}

\section{Open Questions}\label{open-questions}

\textbf{9.1 The formal equivalence of statistical and structural
necessity.} Sidis argues that animate matter is statistically necessary
given time-symmetric microscopic laws. The tholonic model argues that N
emergence is structurally necessary given D-C interaction under the
minimality conditions of
\href{https://github.com/baardev/truevalue/raw/main/docnav/Research/papers/3_minimal-recursive-triadic-framework/3_minimal-recursive-triadic-framework.pdf}{Paper
3}
\href{https://github.com/baardev/truevalue/raw/main/docnav/Research/papers/3_minimal-recursive-triadic-framework/3_minimal-recursive-triadic-framework.pdf}{Mil24c}.
Are these equivalent claims? If yes, the tholonic framework provides a
structural proof of what Sidis could only argue statistically, and the
dual entropy gradients follow as a corollary.

\textbf{9.2 The Landauer bridge.} Landauer's principle states that
erasing one bit of information dissipates at least \(kT \ln 2\) of
energy into the material domain. This is a quantitative connection
between conceptual operations and material entropy. Does it provide a
measurable exchange rate between D-domain entropy decrease and C-domain
entropy increase? If the dual gradients are balanced at the N state, the
Landauer rate may set the equilibrium condition.

\textbf{9.3 The starting-point identification.} The material singularity
and the conceptual void are described as the starting points of the two
trajectories. Are they the same starting point? Is the material Big Bang
and the conceptual void the same N state viewed through the D and C
lenses respectively? If yes, the universe has a single parent N and two
diverging trajectories, exactly the tholonic structure.

\textbf{9.4 Kuiper's theorem as a tholonic statement.} Kuiper's theorem
says the infinite-dimensional unitary group is contractible. Is this a
tholonic statement: that the fully realized N state (the plenum of all
unitary transformations) has the same invariant structure as the parent
N (the point from which they all derived)? Translating Kuiper's theorem
into tholonic grammar may provide the formal bridge between the
mathematical result and the structural claim.

\textbf{9.5 Entropy and the tholonic constants.} The five constants
produced by the tholonic ladder recurrence (\(\pi/4\), \(\varphi\),
\(e\), \(\sqrt{2}\), \(\ln 2\)) all appear in entropy-related
expressions: \(\ln 2\) in Landauer's principle and binary Shannon
entropy, \(e\) in Boltzmann's entropy formula, \(\pi\) in the entropy of
the Gaussian distribution. Is the appearance of the tholonic constants
in entropy formulas a structural consequence of the D-C role assignment,
or coincidental?

\begin{center}\rule{0.5\linewidth}{0.5pt}\end{center}

\section{Conclusion}\label{conclusion}

The Tholonic N-D-C framework, applied at the cosmic scale, assigns the
material domain to the C role (expansive, accumulating, dispersive) and
the conceptual domain to the D role (constraining, organizing,
specifying). This assignment predicts dual entropy gradients: material
entropy increases while conceptual entropy decreases. The two
trajectories begin from equivalent starting points (the material
singularity and the conceptual void are both states of zero
differentiation) and converge toward an equivalent endpoint (the plenum
of all realized N states is informationally equivalent to the void of
all potential N states).

This duality was anticipated by Newton's distinction between passive
mechanical principles and active organizing principles, and by Sidis's
distinction between inanimate (entropy-increasing) and animate
(entropy-reversing) matter. Both investigators were identifying the
interface between the C and D domains, the N state where material and
conceptual entropy gradients meet and locally balance. Neither had the
structural framework to explain why the interface exists or why the two
gradients are opposed.

The tholonic model provides that framework. The dual entropy gradients
are not empirically independent observations. They are the C and D roles
operating at the largest observable scale, with the universe itself as
the N state of their ongoing negotiation.

The void and the plenum are the same entity. They are the parent N and
the fully realized child N of the same recursive tholonic hierarchy,
indistinguishable from outside the recursion because both are states of
complete determination: one in potential, one in actuality. The universe
is the process by which one becomes the other.

\begin{center}\rule{0.5\linewidth}{0.5pt}\end{center}

\section{References}\label{references}

{[}Bol77{]} Boltzmann, L. ``Über die Beziehung zwischen dem zweiten
Hauptsatze der mechanischen Wärmetheorie und der
Wahrscheinlichkeitsrechnung.'' \emph{Wiener Berichte}, 76:373-435, 1877.

{[}Kth65{]} Kuiper, N. H. ``The homotopy type of the unitary group of
Hilbert space.'' \emph{Topology}, 3(1-2):19-30, 1965.

{[}Lan61{]} Landauer, R. ``Irreversibility and heat generation in the
computing process.'' \emph{IBM Journal of Research and Development},
5(3):183-191, 1961.

\href{https://github.com/baardev/truevalue/raw/main/docnav/Research/papers/1_recursive-tholonic-five-constants/1_recursive-tholonic-five-constants.pdf}{Mil24a}
Milton, J. W. ``Emergence of Classical Constants from a Minimal
Recursive Triadic Framework.'' Clarity Coalition Working Paper, 2024.

\href{https://github.com/baardev/truevalue/raw/main/docnav/Research/papers/2_supply-chain-transparency-tvpci/2_supply-chain-transparency-tvpci.pdf}{Mil24b}
Milton, J. W. ``Phase-Resolved Transparency Classification (TVPCI).''
Clarity Coalition Working Paper, 2024.

\href{https://github.com/baardev/truevalue/raw/main/docnav/Research/papers/3_minimal-recursive-triadic-framework/3_minimal-recursive-triadic-framework.pdf}{Mil24c}
Milton, J. W. ``A Minimal Recursive Triadic Framework for Self-Similar
Hierarchical Systems.'' Clarity Coalition Working Paper, 2024.

\href{https://github.com/baardev/truevalue/raw/main/docnav/Research/papers/12_newton-tholonic-framework/12_newton-tholonic-framework.pdf}{Mil26a}
Milton, J. W. ``Newton's Framework Through the Tholonic Lens: Alchemy,
Physics, and the Triadic Structure of Nature.'' Clarity Coalition
Working Paper, 2026.
(\href{https://github.com/baardev/truevalue/raw/main/docnav/Research/papers/12_newton-tholonic-framework/12_newton-tholonic-framework.pdf}{Paper
12 in this series}.)

{[}New17{]} Newton, I. \emph{Opticks: Or, A Treatise of the Reflections,
Refractions, Inflections and Colours of Light}. Second edition. London,
1717. (Query 31.)

{[}Sha48{]} Shannon, C. E. ``A Mathematical Theory of Communication.''
\emph{Bell System Technical Journal}, 27:379-423, 623-656, 1948.

{[}Sid25{]} Sidis, W. J. \emph{The Animate and the Inanimate}. Badger,
1925.

{[}Whe90{]} Wheeler, J. A. ``Information, physics, quantum: the search
for links.'' In \emph{Complexity, Entropy, and the Physics of
Information}, ed.~W. H. Zurek. Addison-Wesley, 1990.

\end{document}
